% !PRINTED 436
\weberpage{436}
But all the roots of (3) are, by § 69:
\weq{5}{j(p\omega), \quad j\left(\frac{\omega + \lambda}{p}\right),}
where $\lambda$ runs through a system of residues with respect to the
modulus $p$. The first of these does not belong to the roots of (4), for
$\omega' = p\omega$ satisfies the equation:
\[
a\omega'^2 + p b \omega' + p^{\varepsilon+2} c = 0,
\]
which is primitive and of discriminant $-mp^2$. If, however, we set
$p\omega' = \omega + \lambda$, we obtain for $\omega'$ the equation:
\weq{6}{a p^2 \omega'^2 + p(b - 2a\lambda)\omega' + (a\lambda^2 - b\lambda + c p^{\varepsilon}) = 0,}
and these equations likewise have the discriminant
$-mp^2$. But there are two among them which are imprimitive with the
divisor $p$, namely those which correspond to the values
\weq{7}{\lambda \equiv 0, \quad a\lambda - b \equiv 0 \pmod{p}}
The first is the root of the form $(a p, b, c p^{\varepsilon-1})$, the
second that of the form
\[
(a p, \; b - 2a\lambda, \; C),
\]
and since $b - 2a\lambda \equiv b \pmod{2a}$ and [by (7)] $\equiv -b \pmod{2p}$,
this form is compounded from $(a, b, c p^{\varepsilon})$ and $(p, -b, a c p^{\varepsilon-1})$,
which belongs to the class $l^{-1}$ reciprocal to $l$. From this it follows:

I. \gesperrt{The equations (3) and (4) have only two roots
 in common with one another, namely the class invariants
$(lk)$ and $(l^{-1}k)$, and if the class $l$
is two-sided, they have only one root
in common.}

If, then, one sets $u = (k)$ in (3), the greatest common
 divisor of (3) and (4) is of the second degree, and if $l$
is two-sided, of the first degree. In the former case $(lk)$ and
$(l^{-1}k)$ are the roots of a quadratic equation whose coefficients
 depend rationally on $u$, and whose form is for the rest
dependent only on the class $l$, not on $k$. To what has thus
been proved we can give the following expression:

II. \gesperrt{We have}
\weq{8}{(lk) + (l^{-1}k) = f_l(k),}
\gesperrt{or, if $l$ is two-sided,}
\weq{9}{(lk) = f_l(k),}
\gesperrt{in which $f_l(k)$ denotes a rational function of $(k)$,
whose rational numerical coefficients depend only
on the class $l$, not on the class $k$.}
